%%%%%%%%%%%%%%%%%%%%%%%%%%%%%%%%%%%%%%
% Cover Letter - Backend Developer .NET SSR (vía ACTIF)
%%%%%%%%%%%%%%%%%%%%%%%%%%%%%%%%%%%%%%

\documentclass[]{cover}
\usepackage{fancyhdr}

\pagestyle{fancy}
\fancyhf{}

\rfoot{Page \thepage \hspace{0pt}}
\thispagestyle{empty}
\renewcommand{\headrulewidth}{0pt}
\begin{document}

%%%%%%%%%%%%%%%%%%%%%%%%%%%%%%%%%%%%%
%     TITLE NAME
%%%%%%%%%%%%%%%%%%%%%%%%%%%%%%%%%%%%%
\namesection{Walter}{Hernandez}{  \href{mailto:walskp@gmail.com}{walskp@gmail.com} | (+54) 11 2382-6514 |  \urlstyle{same}\href{https://www.linkedin.com/in/wal-hernandez-dev}{LinkedIn} | \urlstyle{same}\href{https://github.com/Wal-Hernandez}{GitHub}
}

%%%%%%%%%%%%%%%%%%%%%%%%%%%%%%%%%%%%%
%     MAIN COVER LETTER CONTENT
%%%%%%%%%%%%%%%%%%%%%%%%%%%%%%%%%%%%%

\currentdate{\today}
\lettercontent{Estimado equipo de ACTIF,}

\lettercontent{Me presento para aplicar a la posición Backend Developer .NET SSR que están gestionando. Con experiencia en desarrollo backend con C\#, .NET, Node.js y AWS, me entusiasma esta oportunidad porque combina el trabajo con .NET y servicios cloud en los que tengo experiencia práctica con el entorno SaaS y la integración con inteligencia artificial, áreas en las que quiero seguir creciendo.}

\lettercontent{En mi rol actual en La Nación contribuí a la migración de un sistema legacy .NET 2.2 de newsletters hacia una arquitectura event-driven en Node.js y AWS, desarrollando APIs REST y servicios backend con Lambda, EventBridge y SQS, y aprovisionando infraestructura con AWS CDK. También desplegué servicios mediante pipelines CI/CD en Azure DevOps y realicé code reviews. Mi experiencia se alinea directamente con lo que busca la posición:}

{\raggedright\fontspec[Path = \coverfontpath/raleway/]{Raleway-Medium}\fontsize{11pt}{13pt}\selectfont
\begin{itemize}
    \item \textbf{Backend y APIs:} Diseño APIs REST y modelos de datos en C\#/.NET y Node.js, aplico principios SOLID y patrones de diseño, y mantengo documentación técnica.
    \item \textbf{Servicios backend y cloud:} Desarrollé servicios backend orientados a eventos con EventBridge, SQS y Lambda. Provisioné infraestructura con AWS CDK y configuré monitoreo con CloudWatch.
    \item \textbf{Funcionalidades con IA:} En Eppical integré OpenAI APIs en flujos de procesamiento documental, procesando resultados en servicios C\#/.NET y Node.js.
    \item \textbf{Code reviews y CI/CD:} Realizo code reviews, despliego mediante Azure DevOps y depuro problemas en producción trazando eventos a través de EventBridge y logs de Lambda.
\end{itemize}\par}
\vspace{6pt}

\lettercontent{Me motiva especialmente la combinación de un producto SaaS B2B con alcance en Latinoamérica, una arquitectura moderna sobre AWS y una cultura colaborativa con foco en aprendizaje y feedback. Disfruto tomar ownership técnico de sistemas backend, proponer mejoras basadas en datos y construir software donde la calidad y la mantenibilidad se valoran. Trabajo cómodamente en entornos híbridos o remotos y en equipos multidisciplinarios.}

\lettercontent{Me encantaría conversar sobre cómo mi experiencia en backend, cloud y .NET puede contribuir al equipo de ingeniería de su cliente. Quedo a disposición para ampliar cualquier información.}

\begin{flushright}
\closing{Saludos cordiales,}

\signature{Walter Hernandez}
\end{flushright}
\end{document}
